\begin{table}
\caption{描述性统计：当前候选 DML 样本（2019—2023）}
\label{tab:descriptive_statistics}
\begin{tabular}{llllllll}
\toprule
变量 & 角色 & 样本量 & 均值 & 标准差 & 最小值 & 中位数 & 最大值 \\
\midrule
碳排放强度 & 主结果变量 & 1456 & 14.6691 & 16.9532 & 1.4320 & 9.8553 & 198.5941 \\
碳排放总量 & 稳健性结果变量 & 1456 & 34046.4204 & 40815.1023 & 929.8957 & 23344.5090 & 440958.3002 \\
数字普惠金融指数 & 处理变量 & 1456 & 275.5286 & 31.6054 & 192.3858 & 275.1887 & 373.2215 \\
覆盖广度 & 数字金融分项 & 1456 & 282.5982 & 44.1354 & 179.5993 & 284.8138 & 408.2614 \\
使用深度 & 数字金融分项 & 1456 & 250.4614 & 28.5537 & 165.8475 & 248.9009 & 354.3049 \\
数字化程度 & 数字金融分项 & 1456 & 297.7528 & 16.3626 & 247.1270 & 297.1850 & 342.3680 \\
地区生产总值 & 当前主规格控制变量 & 1456 & 3690.1360 & 5304.6472 & 187.0000 & 2107.0000 & 47219.0000 \\
第二产业占比 & 当前主规格控制变量 & 1456 & 39.3834 & 10.4484 & 10.6800 & 39.6500 & 73.0300 \\
财政支出 & 当前主规格控制变量 & 1456 & 5988317.8915 & 8384544.5681 & 268050.0000 & 4099725.0000 & 96385132.0000 \\
人口规模（候选控制） & 候选控制变量，未锁定 & 1456 & 453.6209 & 356.9442 & 21.0000 & 372.0000 & 3416.0000 \\
\bottomrule
\end{tabular}
\end{table}
